\documentclass[12pt]{article}
%^ Start of document
\usepackage{times}  %Makes text TimesNewRoman
\usepackage{setspace} %Allows you to set single or double space 
\usepackage{biblatex} %Imports biblatex package
\usepackage{graphicx} % Required for inserting images
\usepackage{authblk}
\usepackage[colorlinks=true, urlcolor=blue, linkcolor=blue]{hyperref}
\usepackage{subfig}
\usepackage{float}
\usepackage{tabularx}
\usepackage{multirow}
\usepackage[paperheight=11in,
   paperwidth=8.5in,
   top=20mm,
   bottom=20mm,
   left=40mm,
   right=40mm]{geometry}
\usepackage{moresize}
\usepackage{tikz}
\usepackage{xcolor}
\usepackage{physics}
\usepackage{amssymb}
\usepackage{mathrsfs}
\usepackage{amsmath}
\usepackage{ragged2e}
\usepackage{comment}
\usepackage{xcolor}
\addbibresource{seminar.bib}
\begin{document}
\singlespacing  %Sets single spacing for whole document
\title{Biofilms - McLean}

\author[1]{Zachary Tullier}
\affil[1]{Student\\Physics Department\\ University of Houston - Clear Lake \\Houston, Texas\\U.S}
\date{}
\maketitle



\section{Introduction and Personal Background}
    \textbf{Speaker:} Dr. Robert (Bob) McLean is a renowned microbiologist specializing in biofilms and their behavior in diverse environments, including unique conditions like microgravity. Originally from Canada, Dr. McLean completed his undergraduate studies at the University of Guelph and earned his Ph.D. from the University of Calgary, followed by a postdoctoral fellowship at Guelph. His career includes roles in the food industry, academia, and extensive collaborations with NASA on biofilm research. Inspired by watching John Glenn’s space launch as a child, Dr. McLean eventually realized his dream of contributing to space exploration through scientific research. Dr. McLean has been a professor at Texas State University since 1993, leading groundbreaking research in microbiology. His work focuses on biofilms, microbiological processes, and their implications for natural ecosystems and human-made systems.

\section{Understanding Biofilms: Definitions and Importance}
    Biofilms are complex communities of microorganisms that adhere to surfaces and produce a slimy extracellular matrix for protection. They are found in natural environments such as rivers and lakes and human systems like industrial pipelines and medical devices, biofilms are both beneficial and problematic. Biofilms form when bacteria attach to a surface, multiply, and produce a matrix of sugars, proteins, and DNA that holds the community together. These structures protect bacteria from external threats, including antibiotics and disinfectants, making them highly resistant to treatment. Biofilms play key roles in nutrient cycling and environmental stability, such as in aquatic ecosystems. They can also cause problems, including tooth decay, medical device infections, pipeline corrosion, and contamination of industrial systems.

\section{Pioneering Biofilm Research in Microgravity}
    Dr. McLean’s interest in biofilms in space began with a science fair collaboration with eighth-grade students, exploring whether biofilms could form in low gravity. This initial experiment flew on the Space Shuttle Discovery in 1998, demonstrating that biofilms can grow in microgravity. Dr. McLean’s recent research focuses on biofilm formation aboard the International Space Station (ISS) and the challenges it poses to spacecraft systems. These studies were conducted during SpaceX missions, including SpaceX-21 and SpaceX-29, using custom-designed experimental devices. Understand how biofilms grow, respond to disinfectants, and contribute to corrosion in spacecraft environments. Explore the genetic and functional changes in bacteria exposed to microgravity and space radiation.

\section{Key Findings and Insights}
    Biofilms were found to form successfully in microgravity, attaching to surfaces and developing structural characteristics similar to those on Earth. Bacteria in space were observed to undergo mutations, potentially influenced by the absence of gravity, radiation exposure, or other environmental factors. Biofilms exhibited heightened resistance to disinfectants like silver ions, a commonly used antimicrobial agent aboard the ISS. Resistance mechanisms may include genetic mutations, efflux pumps (systems that expel harmful substances from bacterial cells), and the protective nature of the biofilm matrix. Biofilms have been identified as a major issue for spacecraft life support systems, particularly in water recovery systems aboard the ISS. The slimy nature of biofilms can clog filters, valves, and pipes, leading to failures that could pose serious risks during long-term missions, such as those to Mars. During one experiment, bacterial samples returned from a 117-day spaceflight showed evidence of living bacteria with possible mutations, while similar ground-based samples did not display these changes. Advanced analysis, including genetic sequencing, is being used to identify the nature of these mutations and their implications for bacterial behavior in space.


\section{Broader Implications of Biofilm Research}
    Biofilm-related corrosion and contamination are global challenges, costing billions annually in industries like oil, water, and manufacturing. Insights from space research can inform strategies to mitigate biofilm formation and damage in terrestrial systems. Collaboration with computer scientists has led to the development of machine learning tools to detect and classify corrosion caused by biofilms. DNA and RNA analysis of space-exposed bacteria is providing valuable information on how biofilms adapt to extreme environments.


\section{The Interdisciplinary Nature of Biofilm Studies}
    Dr. McLean’s research has brought together experts in microbiology, genetics, engineering, computer science, and space science. NASA engineers and astronauts played a critical role in designing and executing experiments aboard the ISS, demonstrating the importance of teamwork in addressing complex challenges. Dr. McLean has mentored students and researchers, integrating real-world biofilm research into education and inspiring the next generation of scientists. He emphasizes the importance of interdisciplinary approaches to solving scientific problems.


\section{Challenges and Future Directions}
    What specific factors in microgravity—such as reduced convection, radiation, or magnetic field effects—drive the observed genetic mutations in bacteria? How do these mutations contribute to biofilm resilience and resistance to antimicrobials? Managing biofilms is critical for long-term space missions, where access to resources and emergency interventions are limited. Research aims to predict and mitigate potential system failures caused by biofilms in spacecraft, paving the way for sustainable human exploration of Mars and beyond. Dr. McLean expressed his fascination with the unexpected discoveries and interdisciplinary nature of his research. He highlighted the importance of addressing biofilm-related challenges to ensure the success of future space missions.


\section{Conclusion:}
    Dr. McLean credited his team of students, collaborators, and NASA personnel for their contributions to the success of these experiments. He emphasized the value of teamwork and diverse expertise in solving complex scientific problems.


\end{document}
